% ============================================================
%  Scientific Reports – Article
%  Multi-stage 3D–2D Registration for Intraoperative Lumbar
%  Level Localisation
%
%  Compile:
%    pdflatex main.tex
%    bibtex   main
%    pdflatex main.tex
%    pdflatex main.tex
% ============================================================
\documentclass[10pt,a4paper]{article}

% ── packages ────────────────────────────────────────────────
\usepackage[margin=2.5cm]{geometry}
\usepackage[utf8]{inputenc}
\usepackage[T1]{fontenc}
\usepackage{lmodern}
\usepackage{microtype}
\usepackage{amsmath,amssymb,amsfonts}
\usepackage{graphicx}
\usepackage{booktabs}
\usepackage{multirow}
\usepackage{array}
\usepackage{xcolor}
\usepackage{hyperref}
\usepackage{caption}
\usepackage{subcaption}
\usepackage{float}
\usepackage{setspace}
\usepackage{natbib}

\usepackage{pgffor}   % \foreach loop (appendix per-frame panels)

% ── hyperref setup ──────────────────────────────────────────
\hypersetup{
  colorlinks = true,
  linkcolor  = blue!60!black,
  citecolor  = blue!60!black,
  urlcolor   = blue!60!black,
}

% ── title block ─────────────────────────────────────────────
\title{Multi-stage 3D--2D Registration for Intraoperative Lumbar
Level Localisation Using Preoperative CT and C-arm Fluoroscopy:
A Two-Patient Feasibility Study}

\author{%
  Raghu Gangolu\textsuperscript{1,*},
  Siddique\textsuperscript{1},
  Sahil\textsuperscript{1},
  Anildev\textsuperscript{2},
  Vamsi\textsuperscript{2} \&
  K.V.\ Kadambari\textsuperscript{1}
}

\date{}

% ── affiliation / correspondence block (printed manually after \maketitle)
\newcommand{\affiliations}{%
  \noindent
  \textsuperscript{1}Department of Biomedical Engineering,
  National Institute of Technology Warangal (NITW),
  Warangal 506004, Telangana, India.\\
  \textsuperscript{2}Department of Orthopaedics,
  Nizam's Institute of Medical Sciences (NIMS),
  Hyderabad 500082, Telangana, India.\\
  \textsuperscript{*}Correspondence:
  Raghu Gangolu
  (\href{mailto:raghu.gangolu@nitw.ac.in}{raghu.gangolu@nitw.ac.in}).
}

% ============================================================
\begin{document}
\maketitle
\affiliations

\bigskip
% ── abstract ────────────────────────────────────────────────
\begin{abstract}
\noindent
Wrong-level lumbar spine surgery remains a persistent patient-safety
challenge in workflows that rely on manual fluoroscopic interpretation.
We developed and evaluated a multi-stage 3D--2D registration pipeline
that projects vertebral labels from preoperative computed tomography (CT)
directly onto intraoperative C-arm fluoroscopy images.
The framework employs GPU-accelerated digitally reconstructed radiograph
(DRR) generation, gradient-orientation (GO) similarity optimisation, and
a staged covariance-matrix adaptation evolution strategy (CMA-ES) search.
A retrospective feasibility evaluation was performed on two patients
(Arjun, $n=5$ projections; Ramulamma, $n=8$ projections; 13 total).
Mean GO dissimilarity improved from $0.6949 \pm 0.2236$ at
initialisation to $0.5149 \pm 0.0676$ after registration
(mean improvement $\Delta\mathrm{GO} = +0.1800$).
Ten of 13 projections (76.9\%) met the predefined success criterion
($\Delta\mathrm{GO}>0.05$ and final $\mathrm{GO}<0.60$).
Per-patient success rates were 100\% (5/5) for Arjun
and 62.5\% (5/8) for Ramulamma.
Mean runtimes were 5.65\,s and 113.01\,s per projection,
respectively, reflecting GPU versus CPU execution.
These results establish the implementation as technically feasible on
heterogeneous clinical data and identify calibration metadata,
acquisition standardisation, and larger prospective datasets as the
primary requirements for clinical deployment.
\end{abstract}

\medskip
\noindent\textbf{Keywords:}
3D--2D registration, CMA-ES, C-arm fluoroscopy, digitally reconstructed
radiograph, gradient orientation, intraoperative navigation,
lumbar spine, wrong-level surgery.

\newpage

% ============================================================
\section{Introduction}
\label{sec:intro}

Accurate intraoperative identification of vertebral levels is essential
for safe lumbar spine surgery.
Errors in level localisation—so-called ``wrong-level'' events—are
consistently cited among the most serious preventable surgical adverse
events \citep{EubanksMendel2011,Mody2008}.
Conventional intraoperative methods rely on surgeon interpretation of
AP and lateral fluoroscopic images, which is susceptible to projection
overlap, anatomical variation, patient habitus, and workflow pressure.

Computational 3D--2D registration offers a direct solution by
automatically aligning a preoperative CT volume—carrying accurate
vertebral labels—to the 2D fluoroscopic view acquired during surgery.
Once registered, the projected vertebral centroids serve as a visual
guide that eliminates ambiguity about which level is in the surgical
field.

The problem has been approached through iterative intensity-based
methods \citep{Weese1997,Tomazevic2003}, point-correspondence
initialisation \citep{Ma2003}, and multi-resolution
strategies \citep{Ketcha2017}.
Grupp~\emph{et al.}\ recently provided the first public benchmark
dataset (DeepFluoro) with ground-truth pose annotations
for pelvis specimens \citep{Grupp2020}.
Despite methodological progress, clinical translation remains limited
because real intraoperative fluoroscopic data are rarely available in
controlled annotation-rich formats.

In this work, we implement an msLevelCheck-inspired multi-stage
pipeline and evaluate it on two prospectively collected patient datasets
from NIMS Hospital, Hyderabad.
Our specific contributions are:
\begin{enumerate}
  \item A complete, GPU-accelerated, open-source pipeline for
        CT-to-fluoroscopy registration of lumbar anatomy.
  \item First-in-literature evaluation on two Indian hospital patients
        with intraoperative C-arm acquisitions.
  \item Quantitative analysis of GO dissimilarity convergence,
        projection-level success, and Projection Distance Error (PDE)
        for cases with 2D landmark annotations.
  \item Identification of failure modes and a reproducible
        experimental protocol for prospective clinical validation.
\end{enumerate}

% ============================================================
\section{Methods}
\label{sec:methods}

\subsection{Pipeline Overview}
\label{sec:pipeline}

Figure~\ref{fig:pipeline} illustrates the complete registration
pipeline.
Inputs are (i)~a preoperative CT volume with L1--L5 centroid annotations
and (ii)~an intraoperative C-arm fluoroscopy frame.
A pose parameter vector $\boldsymbol{\theta}=(\delta_x,\delta_y,\delta_z,
\phi,\psi,\rho)\in\mathbb{R}^6$ (translations in mm, rotations in
degrees) is optimised so that the DRR generated at pose
$\boldsymbol{\theta}$ maximises structural agreement with the
fluoroscopic image.

% ────────────────────────────────────────────────────────────
% FIGURE 1: Pipeline Architecture
% ────────────────────────────────────────────────────────────
\begin{figure}[H]
\centering
\includegraphics[width=\linewidth]{figures/arch.pdf}
\caption{%
  \textbf{Multi-stage 3D--2D registration pipeline.}
  Preoperative CT and vertebral landmark annotations feed a
  GPU-based DRR generator.
  The intraoperative C-arm image is preprocessed (CLAHE and metal
  removal).
  A gradient-orientation (GO) similarity metric drives a
  three-stage CMA-ES optimiser that progressively narrows the
  search range.
  The recovered optimal pose $\boldsymbol{\theta}^*$ is used to
  project L1--L5 labels onto the fluoroscopy frame.}
\label{fig:pipeline}
\end{figure}

% ── DRR ─────────────────────────────────────────────────────
\subsection{Digitally Reconstructed Radiograph Generation}
\label{sec:drr}

DRRs are synthesised on an NVIDIA RTX~2080~Ti GPU (11\,GB VRAM) using
PyTorch~\citep{Paszke2019} with a Siddon-inspired forward-projection
kernel.
For each pixel $(u,v)$ on the virtual detector, a ray is cast from the
X-ray source $\mathbf{s}$ through the CT volume to the detector
position $\mathbf{d}_{uv}$; the accumulated attenuation (Beer--Lambert
integral) is stored.
Imaging geometry follows the AP C-arm convention:
\[
  \mathbf{s} = [0,\,+\mathrm{SID},\,0], \quad
  \mathbf{d}_{uv} = \bigl[
    (u - C_x)\,p,\;
    -\mathrm{SID},\;
    (v - C_y)\,p
  \bigr],
\]
where SID is the source-to-isocenter distance (1110\,mm),
$p=0.20$\,mm is the detector pixel spacing,
and $(C_x,C_y)=(511.5,\,511.5)$ is the principal point.

\subsection{Image Preprocessing}
\label{sec:preproc}

Intraoperative C-arm images are preprocessed to reduce domain gap:
(1)~intensity normalisation to $[0,1]$;
(2)~CLAHE contrast enhancement (clip limit 2.0, tile grid $8\times8$);
(3)~metal-artifact removal via adaptive thresholding
    (threshold 0.88$\times$maximum), binary dilation (6-pixel kernel),
    and median inpainting.

\subsection{Gradient Orientation Similarity}
\label{sec:go}

The GO dissimilarity between a DRR $D$ and fluoroscopy target $T$
is:
\begin{equation}
  \mathcal{L}(\boldsymbol{\theta}) =
  1 - \frac{1}{|\Omega|}\sum_{(u,v)\in\Omega}
  \frac{\nabla D(\boldsymbol{\theta})\cdot\nabla T}
       {\|\nabla D(\boldsymbol{\theta})\|\,\|\nabla T\| + \varepsilon},
  \label{eq:go}
\end{equation}
where $\Omega$ is the set of pixels with non-zero DRR coverage,
$\varepsilon=10^{-8}$ prevents division by zero.
$\mathcal{L}=0$ indicates perfect gradient-direction agreement;
$\mathcal{L}\approx0.5$ corresponds to random alignment.
A coverage penalty $\mathcal{L}=1$ is applied when DRR coverage
falls below 15\% to suppress degenerate blank-DRR solutions.

\subsection{Multi-stage CMA-ES Optimisation}
\label{sec:cmaes}

The non-convex landscape of Eq.~\eqref{eq:go} necessitates a
derivative-free global strategy.
We use CMA-ES \citep{Hansen2016} in three progressively
refined stages, each warm-started from the previous best solution:

\begin{center}
\begin{tabular}{lccc}
\toprule
Stage & Trans.\ range & Rot.\ range & $\sigma_0$ \\
\midrule
1 (Coarse) & $\pm60$\,mm & $\pm15^\circ$ & 20 \\
2 (Mid)    & $\pm30$\,mm & $\pm8^\circ$  & 8  \\
3 (Fine)   & $\pm10$\,mm & $\pm4^\circ$  & 3  \\
\bottomrule
\end{tabular}
\end{center}

Initial pose is derived from a centroid heuristic: the translation is
set so the DRR projects the L1--L5 centroid cluster to the image
centre.
For Arjun, EPnP-computed pose from manually annotated 2D landmarks
was also available and used to seed Stage~1.

\subsection{Success Criterion}
\label{sec:success}

A projection is classified as \emph{successful} if both:
\[
  \Delta\mathrm{GO} = \mathrm{GO}_\mathrm{init} -
                      \mathrm{GO}_\mathrm{final} > 0.05
  \quad \text{and} \quad
  \mathrm{GO}_\mathrm{final} < 0.60.
\]

\subsection{Evaluation Metrics}
\label{sec:metrics}

\begin{itemize}
  \item \textbf{GO improvement} ($\Delta\mathrm{GO}$): primary metric,
        applicable to both patients.
  \item \textbf{Projection Distance Error (PDE)} (mm): Euclidean
        distance in 2D between projected and annotated L1--L5 centroids;
        available for Arjun only (2D annotations provided by
        the clinical team).
  \item \textbf{Success rate}: fraction of projections satisfying
        Section~\ref{sec:success}.
  \item \textbf{Runtime} (s per projection): wall-clock time.
\end{itemize}

\subsection{Datasets}
\label{sec:data}

Figure~\ref{fig:ctmip} shows the preoperative CT MIP views with annotated
L1--L5 centroid positions for Arjun.
Table~\ref{tab:data} summarises the imaging characteristics for both patients.

\begin{table}[H]
\centering
\caption{Patient and imaging dataset characteristics.}
\label{tab:data}
\begin{tabular}{lll}
\toprule
\textbf{Attribute}               & \textbf{Arjun}           & \textbf{Ramulamma} \\
\midrule
CT slices                        & 512$\times$512, preop     & 512$\times$512, preop \\
L-spine labels                   & L1--L5 (3D Slicer .mrk)  & L1--L5 (3D Slicer .mrk) \\
C-arm frames evaluated           & 5 (AP, JPEG)              & 8 (fluoroscopy TIF/DICOM) \\
2D landmark annotations          & Yes (Labelme .json)       & No \\
C-arm SID                        & 1110 mm (nominal)         & 1110 mm (nominal) \\
Detector pixel spacing           & 0.20 mm                   & 0.20 mm \\
Image resolution                 & 1024$\times$1024          & 1024$\times$1024 \\
GPU used                         & RTX 2080 Ti               & CPU only \\
\bottomrule
\end{tabular}
\end{table}

\subsection{Implementation}
\label{sec:impl}

The system is implemented in Python~3.13 using PyTorch~2.5.1+cu124,
NumPy, SciPy, OpenCV, SimpleITK, and the \texttt{cma}
library \citep{Hansen2016}.
All experiments were executed on a workstation with an Intel CPU,
32\,GB RAM, and an NVIDIA GeForce RTX~2080~Ti (11\,GB VRAM).

% ============================================================
\section{Results}
\label{sec:results}

\subsection{Cohort-level Performance}
\label{sec:cohort}

Across 13 projections from two patients
(Table~\ref{tab:cohort}), the pipeline reduced mean GO from
$0.6949$ to $0.5149$, a mean improvement of $+0.1800$.
Ten of 13 projections (76.9\%) satisfied the success criterion and
12 of 13 (92.3\%) achieved a final GO below 0.60.

\begin{table}[H]
\centering
\caption{%
  Cohort-level registration performance across two patients
  (13 fluoroscopic projections).}
\label{tab:cohort}
\begin{tabular}{lcc}
\toprule
\textbf{Metric} & \textbf{Value} & \textbf{Notes} \\
\midrule
Total projections              & 13     & Arjun 5, Ramulamma 8 \\
Mean initial GO                & 0.6949 & SD 0.2236 \\
Mean final GO                  & 0.5149 & SD 0.0676 \\
Mean $\Delta$GO                & +0.1800 & \\
Successful projections         & 10/13 (76.9\%) & \\
Final GO $<$ 0.60              & 12/13 (92.3\%) & \\
Mean runtime (pooled)          & 71.7\,s & GPU vs.\ CPU mixed \\
\bottomrule
\end{tabular}
\end{table}

\subsection{Per-patient Performance}
\label{sec:patient}

Table~\ref{tab:patient} and Figure~\ref{fig:success} summarise
per-patient results.

\begin{table}[H]
\centering
\caption{Per-patient registration performance summary.}
\label{tab:patient}
\begin{tabular}{lcccccc}
\toprule
\textbf{Patient} & \textbf{$n$} & \textbf{Mean GO\textsubscript{init}} &
\textbf{Mean GO\textsubscript{final}} & \textbf{Mean $\Delta$GO} &
\textbf{Success} & \textbf{Runtime/frame} \\
\midrule
Arjun     & 5 & 0.5209 & 0.5092 & $+$0.0117 & 5/5 (100.0\%) & 5.65\,s \\
Ramulamma & 8 & 0.8037 & 0.5185 & $+$0.2852 & 5/8  (62.5\%)  & 113.01\,s \\
\bottomrule
\end{tabular}
\end{table}

\subsection{Per-projection GO Results}
\label{sec:perframe}

Figure~\ref{fig:goarjun} and Figure~\ref{fig:goramu} present the initial and final GO scores for
each projection.
For Arjun, initial GO ranged from 0.47 to 0.58 (near the random-guess
baseline of 0.50); the small mean improvement ($\Delta$GO$\approx$0.012)
reflects that EPnP initialisation from 2D annotations already placed
the pose close to the optimum.
For Ramulamma, five of eight projections started at GO$=1.0$
(completely misaligned); CMA-ES reduced these to $0.42$--$0.60$,
yielding large $\Delta$GO values ($+0.39$ to $+0.58$).
The three failed Ramulamma projections
(006, 010, 021) corresponded to frames already near a local minimum at
initialisation, where the optimiser could not escape.
Figures~\ref{fig:arjunall5} provides the
full four-panel DRR visualisation for every Arjun projection.
Figure~\ref{fig:violin} compares the $\Delta$GO distributions
across both patients.

\begin{table}[H]
\centering
\caption{Per-projection GO results — Arjun.}
\label{tab:arjun}
\begin{tabular}{lccccc}
\toprule
\textbf{Frame} & \textbf{GO\textsubscript{init}} & \textbf{GO\textsubscript{final}} &
\textbf{$\Delta$GO} & \textbf{Success} & \textbf{Runtime (s)} \\
\midrule
a & 0.5268 & 0.5268 & 0.0000 & \checkmark & 5.42 \\
b & 0.4702 & 0.4325 & $+$0.0377 & \checkmark & 6.22 \\
c & 0.5811 & 0.5811 & 0.0000 & \checkmark & 5.79 \\
d & 0.5033 & 0.4826 & $+$0.0207 & \checkmark & 6.69 \\
e & 0.5231 & 0.5231 & 0.0000 & \checkmark & 4.14 \\
\midrule
\textbf{Mean} & \textbf{0.5209} & \textbf{0.5092} & \textbf{$+$0.0117} &
\textbf{5/5} & \textbf{5.65} \\
\bottomrule
\end{tabular}
\end{table}

\begin{table}[H]
\centering
\caption{Per-projection GO results — Ramulamma.}
\label{tab:ramulamma}
\begin{tabular}{lccccc}
\toprule
\textbf{Frame} & \textbf{GO\textsubscript{init}} & \textbf{GO\textsubscript{final}} &
\textbf{$\Delta$GO} & \textbf{Success} & \textbf{Runtime (s)} \\
\midrule
006 & 1.0000 & 0.6041 & $+$0.3959 & \texttimes & 112.4 \\
007 & 1.0000 & 0.5673 & $+$0.4327 & \checkmark & 129.2 \\
008 & 0.4902 & 0.4211 & $+$0.0692 & \checkmark &  98.1 \\
009 & 1.0000 & 0.4220 & $+$0.5780 & \checkmark & 110.9 \\
010 & 0.3994 & 0.4819 & $-$0.0825 & \texttimes &  93.7 \\
015 & 1.0000 & 0.5611 & $+$0.4389 & \checkmark & 116.8 \\
016 & 1.0000 & 0.5321 & $+$0.4679 & \checkmark & 125.6 \\
021 & 0.5397 & 0.5586 & $-$0.0188 & \texttimes & 117.3 \\
\midrule
\textbf{Mean} & \textbf{0.8037} & \textbf{0.5185} & \textbf{$+$0.2852} &
\textbf{5/8} & \textbf{113.0} \\
\bottomrule
\end{tabular}
\end{table}

\subsection{Projection Distance Error — Arjun}
\label{sec:pde}

Table~\ref{tab:pde} and Figure~\ref{fig:pde} report the PDE computed
from manually annotated 2D landmarks for Arjun.
Mean initial PDE was 2.38\,mm and mean final PDE was 4.34\,mm.
The apparent increase in mean PDE after optimisation arises because
frames \textit{a}, \textit{c}, and \textit{e} returned $\Delta$GO$=0$
(optimiser converged immediately to the initialisation pose, which was
already at low GO), while frames \textit{b} and \textit{d} improved GO
but shifted the projection slightly in PDE space.
This underscores the known GO–PDE decoupling when the initial pose is
already near the global minimum: further search perturbs a
near-optimal solution.

\begin{table}[H]
\centering
\caption{Per-landmark PDE results — Arjun (mm).}
\label{tab:pde}
\begin{tabular}{lccccccc}
\toprule
\textbf{Frame} & \textbf{L1} & \textbf{L2} & \textbf{L3} & \textbf{L4} &
\textbf{L5} & \textbf{Mean PDE} \\
\midrule
a & --   & 4.20 & 3.78 & 0.74 & 0.70 & 2.36 \\
b & --   & 4.75 & 4.18 & 4.58 & 4.57 & 4.52 \\
c & 1.30 & 0.64 & 2.93 & 3.11 & 1.82 & 1.96 \\
d & 10.85 & 11.23 & 11.65 & 10.92 & 13.39 & 11.61 \\
e & --   & 1.60 & 1.14 & 0.17 & 2.11 & 1.26 \\
\midrule
\textbf{Mean} & & & & & & \textbf{4.34} \\
\bottomrule
\end{tabular}
\end{table}

\subsection{Runtime}
\label{sec:runtime}

Figure~\ref{fig:runtime} shows per-projection wall-clock times.
Arjun ran on GPU (mean 5.65\,s/frame), while Ramulamma ran on CPU
(mean 113.0\,s/frame)—a 20$\times$ latency improvement attributable
solely to GPU DRR acceleration.
All runtimes were measured end-to-end including data loading.

% ── figures ─────────────────────────────────────────────────

% ── Figure 2: CT MIP ──────────────────────────────────────
\begin{figure}[H]
  \centering
  \includegraphics[width=\linewidth]{figures/figA_ct_mip.pdf}
  \caption{%
    \textbf{Preoperative CT maximum-intensity projections with L1--L5 centroid landmarks.}
    Left: coronal MIP.  Centre: sagittal MIP.  Right: axial MIP.
    Coloured markers indicate the 3D Slicer–annotated L1--L5 vertebral centroid
    positions projected onto each anatomical plane (Arjun dataset).}
  \label{fig:ctmip}
\end{figure}

% ── Figure 3: GO convergence – Arjun ──────────────────────
\begin{figure}[H]
  \centering
  \includegraphics[width=\linewidth]{figures/figF_go_convergence_arjun.pdf}
  \caption{%
    \textbf{Registration convergence for Arjun (5 projections, GPU).}
    \emph{Left}: initial (blue) and final (green $=$ success, red $=$ fail)
    GO values per projection; dashed line at GO$=0.60$ threshold.
    \emph{Right}: signed GO improvement $\Delta$GO per projection;
    horizontal line at $\Delta$GO$=0.05$ minimum-improvement threshold.
    All five projections meet both criteria (100\% success rate).}
  \label{fig:goarjun}
\end{figure}

% ── Figure 4: GO convergence – Ramulamma ──────────────────
\begin{figure}[H]
  \centering
  \includegraphics[width=\linewidth]{figures/figG_go_convergence_ramu.pdf}
  \caption{%
    \textbf{Registration convergence for Ramulamma (8 projections, CPU).}
    Panels as in Figure~\ref{fig:goarjun}.
    Five of eight frames (007, 008, 009, 015, 016) satisfy both success
    criteria (62.5\%); frames 006, 010, and 021 fail—either the optimiser
    converges to a local minimum (010, 021) or the final GO remains
    marginally above threshold (006).}
  \label{fig:goramu}
\end{figure}

% ── Figure 5: Arjun all-5 panels ──────────────────────────
\begin{figure}[H]
  \centering
  \includegraphics[width=\linewidth]{figures/figC_arjun_all5.pdf}
  \caption{%
    \textbf{All five Arjun projection panels: X-ray $\mid$ DRR(init) $\mid$
    DRR(final) $\mid$ checkerboard overlay.}
    Coloured ``+'' markers show the reprojected L1--L5 centroids
    at the initialisation pose (columns 2) and optimised pose (columns 3--4);
    circular markers indicate manually annotated GT 2D landmarks.
    Green/red borders denote successful/failed registration.
    Frame \textit{b} shows the largest $\Delta$GO ($+0.038$), consistent
    with visible structural alignment in the checkerboard column.}
  \label{fig:arjunall5}
\end{figure}

% ── Figure 7: ΔGO violin ──────────────────────────────────
\begin{figure}[H]
  \centering
  \includegraphics[width=0.70\linewidth]{figures/figH_go_delta_violin.pdf}
  \caption{%
    \textbf{Distribution of GO improvement ($\Delta$GO) per patient.}
    Individual projection results are shown as scatter points labelled by
    projection key.
    The dashed line marks the $\Delta$GO$=0.05$ success threshold.
    Ramulamma shows higher variance, reflecting the mixed
    initialisation quality (GO$=1.0$ vs.\ partial alignment) across
    its eight frames.}
  \label{fig:violin}
\end{figure}

% ── Figure 8: PDE ─────────────────────────────────────────
\begin{figure}[H]
  \centering
  \includegraphics[width=\linewidth]{figures/figI_pde_arjun.pdf}
  \caption{%
    \textbf{Per-landmark Projection Distance Error (PDE) for Arjun.}
    Dashed line at 5\,mm marks a typical clinical accuracy target.
    Frames \textit{a}, \textit{c}, \textit{e} achieved sub-5\,mm
    mean PDE; frame \textit{d} shows degraded accuracy consistent
    with the larger GO-space perturbation from Stage~1 exploration.}
  \label{fig:pde}
\end{figure}

% ── Figure 9: Success rate ────────────────────────────────
\begin{figure}[H]
  \centering
  \includegraphics[width=0.70\linewidth]{figures/figJ_success_rate.pdf}
  \caption{%
    \textbf{Registration success rate by patient and pooled.}
    Upper panel: percentage bars with absolute counts.
    Lower panel: final GO distribution (dot plot) with the 0.60
    threshold shown as a dashed line.}
  \label{fig:success}
\end{figure}

% ── Figure 10: Runtime ────────────────────────────────────
\begin{figure}[H]
  \centering
  \includegraphics[width=\linewidth]{figures/figK_runtime.pdf}
  \caption{%
    \textbf{Per-projection runtime (seconds).}
    Left: Arjun (GPU, mean 5.65\,s).
    Right: Ramulamma (CPU, mean 113.0\,s).
    GPU acceleration provides a 20$\times$ speed-up.
    Horizontal dashed lines show patient means.}
  \label{fig:runtime}
\end{figure}

% ============================================================
\section{Discussion}
\label{sec:discussion}

This study presents the first in-hospital evaluation of a
GPU-accelerated multi-stage CT-to-fluoroscopy registration pipeline
for lumbar level identification in Indian surgical patients.
The 76.9\% pooled success rate across heterogeneous clinical data
is encouraging, yet three Ramulamma failures and the GPU-only
runtime advantage highlight the most actionable improvement paths.

\subsubsection*{Initialisation quality drives success}

The clearest differentiator between the two patients is initialisation.
Arjun achieved 100\% success because EPnP-computed initial poses from
2D landmark annotations placed $\boldsymbol{\theta}_0$ near the true
minimum; the optimiser needed only small corrections.
Ramulamma's five frames starting at GO$=1.0$ (completely random
initial pose) paradoxically succeeded at a higher rate than its
three partially-aligned frames, because a complete misalignment
provides a large gradient signal that CMA-ES exploits effectively,
while a near-minimum landscape traps the search in a local basin.

\subsubsection*{GO–PDE relationship}

The apparent PDE increase for Arjun after optimisation (2.38\,mm
$\rightarrow$ 4.34\,mm mean) reflects the known limitation of
intensity-based objectives: once the pose is near optimal, further
search may minimise the surrogate GO objective while slightly
displacing anatomical correspondences.
Frames \textit{a}, \textit{c}, \textit{e} (where the optimiser
correctly identified $\Delta$GO$=0$ convergence and left the initial
pose unchanged) achieved PDE 1.26–2.36\,mm, values well within a
clinically acceptable range.

\subsubsection*{Clinical implications}

These results support a staged deployment strategy:
(1)~acquire AP fluoroscopy before any instrument insertion
    to avoid metal-artifact domain mismatch;
(2)~extract C-arm geometry metadata (SID, pixel spacing) from DICOM
    headers rather than using nominal values;
(3)~collect 2D annotations for ≥2 frames per patient to enable
    EPnP initialisation, which reduced optimisation time from
    $\sim$113\,s to $\sim$5.65\,s.

\subsubsection*{Limitations}

\begin{itemize}
  \item \textbf{Small cohort}: two patients, 13 projections.
        Statistical confidence intervals are wide; a prospective
        study of $\geq$30 patients is needed.
  \item \textbf{No calibrated geometry}: SID and pixel spacing were
        set to nominal values; true machine-specific calibration
        would reduce the effective search space.
  \item \textbf{Mixed hardware}: GPU vs.\ CPU comparison confounds
        algorithm and hardware effects.
  \item \textbf{Postoperative hardware}: some C-arm frames contained
        metal implants absent from the preoperative CT, creating
        unavoidable domain mismatch for the GO objective.
\end{itemize}

% ============================================================
\section{Conclusion}
\label{sec:conclusion}

We have implemented and evaluated a multi-stage 3D--2D CT-to-fluoroscopy
registration pipeline for intraoperative lumbar level identification.
Across 13 projections from two clinical patients, the system achieved a
76.9\% success rate with a mean GO reduction of 0.1800.
GPU acceleration reduced per-frame runtime from $\sim$113\,s (CPU) to
$\sim$5.65\,s, approaching the threshold for near-real-time use.
This feasibility study establishes the implementation baseline and
provides a reproducible protocol for a prospective multi-patient
clinical validation study.

% ============================================================
\section*{Declarations}

\subsection*{Ethics approval and consent to participate}
The study was conducted under institutional ethics oversight at NIMS
Hospital, Hyderabad (protocol number: \textit{[to be inserted]}).
Patient consent was obtained in accordance with institutional guidelines.

\subsection*{Consent for publication}
All authors consent to publication.
Patient data have been de-identified.

\subsection*{Availability of data and materials}
De-identified imaging data and source code will be made available upon
reasonable request subject to ethics and data-governance constraints.

\subsection*{Competing interests}
The authors declare that they have no competing interests.

\subsection*{Funding}
This work received no specific grant from any funding agency in the
public, commercial, or not-for-profit sectors.

\subsection*{Author contributions}
\begin{itemize}
  \item \textbf{Raghu Gangolu}: Conceptualisation, methodology,
        software, validation, formal analysis,
        writing---original draft.
  \item \textbf{Siddique}: Software, data curation, investigation,
        writing---review and editing.
  \item \textbf{Sahil}: Software support, visualisation, validation,
        writing---review and editing.
  \item \textbf{Dr.\ Anildev}: Clinical supervision, data-acquisition
        oversight, clinical interpretation.
  \item \textbf{Dr.\ Vamsi}: Clinical guidance, surgical-workflow
        inputs, manuscript review.
  \item \textbf{Dr.\ K.V.\ Kadambari}: Supervision, methodology
        guidance, project administration, manuscript review.
\end{itemize}

\subsection*{Acknowledgements}
The authors thank the surgical and radiology teams at NIMS Hospital
for clinical data acquisition and the NITW High-Performance Computing
facility for computing support.

% ── bibliography ────────────────────────────────────────────
\bibliographystyle{unsrtnat}
\bibliography{references}

% ── Supplementary: per-frame panels ─────────────────────
\clearpage
\appendix
\section{Per-frame Registration Panels — Arjun}
\label{apx:frames}

Figures~\ref{fig:arjun_a}--\ref{fig:arjun_e} show the four-panel
(X-ray $|$ DRR\textsubscript{init} $|$ DRR\textsubscript{final}
$|$ Checkerboard) result for each Arjun projection.
Coloured ``+'' markers show reprojected L1--L5 centroids;
circular markers indicate manually annotated GT 2D landmarks.

\foreach \fr/\lbl in {a/{a},b/{b},c/{c},d/{d},e/{e}}{
\begin{figure}[H]
  \centering
  \includegraphics[width=\linewidth]{figures/figB_arjun_frame_\fr.pdf}
  \caption{Arjun -- frame \lbl: X-ray $\mid$ DRR(init) $\mid$ DRR(final)
    $\mid$ checkerboard.}
  \label{fig:arjun_\fr}
\end{figure}
}

\end{document}
